%!TeX root=../main.tex

\chapter{Inspirações}
\label{cap:inspiracoes}

\section{Isometria}
\label{sec:isometria}

A projeção isométrica é um caso particular da projeção ortogonal que permite a representação
gráfica tridimensional em superfícies bidimensionais. Matematicamente, a isometria é
caracterizada pelo alinhamento dos eixos coordenados (\textbf{X}, \textbf{Y} e \textbf{Z}) em relação ao plano de
projeção, de modo que os ângulos formados entre a projeção de cada par de eixos sejam
exatamente iguais a 120$\degree$. Essa propriedade garante que a escala de medida permaneça
uniforme ao longo de todos os eixos, eliminando os pontos de fuga característicos da
perspectiva cônica, que é a forma como o olho humano enxerga a realidade.


\section{\textit{Low Poly}}
\label{sec:low_poly}


\section{\textit{Fast-Paced}}
\label{sec:fast_paced}
